% ============================================================================
%  APÉNDICE B — Tablas extensas
%
%  Dos vías para una tabla que no cabe en una página:
%
%  1. \aatablelong[n]{csv}{caption}{label} en el CUERPO del artículo.
%     Envuelve en \longtab[n]{...}, que conserva la numeración correlativa y
%     coloca la tabla automáticamente después de los apéndices. Es la vía que
%     A&A recomienda, porque \longtable* solo funciona en onecolumn.
%
%  2. longtable directo dentro de un apéndice en \onecolumn, como abajo.
%
%  Rotadas: \begin{landscape} (paquete lscape) alrededor del longtable, o el
%  entorno sidewaystable* para una única página rotada.
%
%  Las tablas grandes de datos primarios NO van en el artículo: se archivan en
%  el CDS en ASCII con su ReadMe, y en el texto se describe el contenido de
%  cada columna. Ver la skill aa-paper, references/tables.md.
% ============================================================================

\onecolumn
\section{Supplementary tables}
\label{app:tables}

Texto introductorio de las tablas.

% \aatablelong[2]{tables/catalogue.csv}{Full catalogue.}{tab:appcatalogue}

\FloatBarrier
\twocolumn
